\chapter[\emj{🌁}~Puntos de Articulación y Puentes (Tarjan)]{\emj{🌁}~Puntos de Articulación y Puentes (Tarjan)}

\begin{dsaquees}

Una red de ciudades conectadas por carreteras 🛣️. Algunas conexiones
son tan críticas que, si se rompen, dejan zonas incomunicadas:

\begin{itemize}
\item Puente (bridge): una CARRETERA que, si la quitas, parte el mapa en
      pedazos aislados.
\item Punto de articulación (cut vertex): una CIUDAD que, si desaparece
      (con todas sus carreteras), desconecta el mapa.
\end{itemize}

Son los "puntos débiles" de una red. Tarjan los encuentra en un solo
recorrido DFS midiendo, para cada nodo, qué tan "atrás" puede regresar
sin usar la arista por la que llegó.

\end{dsaquees}

\begin{dsaotraforma}

En un grafo no dirigido conexo:
\begin{itemize}
\item Un \textbf{puente} es una arista cuya eliminación aumenta el número
      de componentes conexas.
\item Un \textbf{punto de articulación} es un vértice cuya eliminación
      (y la de sus aristas) hace lo mismo.
\end{itemize}

El algoritmo de Tarjan usa DFS con dos marcas por nodo:
\verb|disc[u]| (tiempo de descubrimiento) y \verb|low[u]| (el menor
\verb|disc| alcanzable desde $u$ usando aristas del árbol DFS y a lo
sumo una arista "de retorno"). Para una arista de árbol $u \to v$:
\begin{itemize}
\item Es \textbf{puente} si \verb|low[v] > disc[u]| (v no puede volver a
      u ni a un ancestro sin esa arista).
\item $u$ es \textbf{articulación} si \verb|low[v] >= disc[u]| (con
      manejo especial de la raíz: lo es si tiene 2+ hijos en el DFS).
\end{itemize}

\end{dsaotraforma}

\begin{dsadiagrama}
Grafo:
   0 ─── 1 ─── 2
         │     │
         3 ─── 4 ─── 5

disc = tiempo de descubrimiento, low = mínimo alcanzable

Arista 4-5: al quitarla, el 5 queda aislado -> PUENTE
Arista 0-1: al quitarla, el 0 queda aislado -> PUENTE
Vertice 1: al quitarlo, el 0 se desconecta -> ARTICULACIÓN
Vertice 4: al quitarlo, el 5 se desconecta -> ARTICULACIÓN

El ciclo 1-2-4-3-1 NO tiene puentes (hay ruta alterna).
\end{dsadiagrama}

\begin{dsacomofunciona}
\begin{verbatim}
disc = orden de descubrimiento,  low = mínimo alcanzable

Grafo:  0 — 1 — 2 — 0   (triángulo)  y  2 — 3   (cola)

DFS desde 0:
  nodo │ disc │ low │ cómo se calcula low
  ─────┼──────┼─────┼──────────────────────────────────
   0   │  0   │  0  │
   1   │  1   │  0  │ back-edge 1→0 (disc 0) baja low a 0
   2   │  2   │  0  │ back-edge 2→0 (disc 0) baja low a 0
   3   │  3   │  3  │ hoja, sin retorno

PUENTE:  arista (u,v) es puente si  low[v] > disc[u]
  arista 2—3:  low[3]=3 > disc[2]=2  →  ¡PUENTE! (al quitarla, 3 aislado)
  arista 0—1:  low[1]=0 ≤ disc[0]=0  →  no (el triángulo tiene ruta alterna)

ARTICULACIÓN:  u es corte si algún hijo v cumple  low[v] ≥ disc[u]
  nodo 2:  low[3]=3 ≥ disc[2]=2  →  2 es punto de articulación
\end{verbatim}
\end{dsacomofunciona}

\begin{lstlisting}
CÓMO FUNCIONA (Tarjan, un solo DFS)
1. disc[u] = low[u] = timer++ al visitar u.
2. Por cada vecino v:
   - Si v es el padre: ignora.
   - Si v ya visitado: low[u] = min(low[u], disc[v])  (arista de retorno)
   - Si no: DFS(v); low[u] = min(low[u], low[v]);
       * low[v] > disc[u]  -> arista (u,v) es PUENTE
       * low[v] >= disc[u] -> u es ARTICULACIÓN (raíz: 2+ hijos)

PSEUDOCÓDIGO (puentes)
──────────────────────────────────────────────────────────────

función dfs(u, padre):
    disc[u] = low[u] = timer++
    PARA cada vecino v de u:
        SI v == padre: continuar
        SI disc[v] != -1:                    // ya visitado
            low[u] = min(low[u], disc[v])
        SINO:
            dfs(v, u)
            low[u] = min(low[u], low[v])
            SI low[v] > disc[u]:
                registrar (u, v) como PUENTE

CÓDIGO C++ (puentes y articulaciones)
──────────────────────────────────────────────────────────────

#include <vector>
using namespace std;

vector<int> disc, low;
vector<bool> esArtic;
vector<pair<int,int>> puentes;
int timer_ = 0;

void dfs(int u, int parent, vector<vector<int>>& adj) {
    disc[u] = low[u] = timer_++;
    int hijos = 0;
    for (int v : adj[u]) {
        if (v == parent) continue;
        if (disc[v] != -1) {
            low[u] = min(low[u], disc[v]);     // arista de retorno
        } else {
            hijos++;
            dfs(v, u, adj);
            low[u] = min(low[u], low[v]);
            if (low[v] > disc[u])
                puentes.push_back({u, v});      // PUENTE
            if (parent != -1 && low[v] >= disc[u])
                esArtic[u] = true;              // ARTICULACIÓN
        }
    }
    if (parent == -1 && hijos > 1)
        esArtic[u] = true;                      // raíz con 2+ hijos
}

COMPLEJIDAD (Big-O)
──────────────────────────────────────────────────────────────

  Recurso     │ Costo
  ────────────┼─────────────
  Tiempo      │ O(V + E)
  Espacio     │ O(V)

* Un solo DFS calcula puentes Y articulaciones a la vez.
\end{lstlisting}

\begin{dsaentrevistas}

✅ "Critical Connections in a Network" (LeetCode 1192): es exactamente
   encontrar todos los puentes. Respuesta directa con Tarjan.

💡 disc vs low: \verb|disc| es "cuándo te descubrí"; \verb|low| es "lo más
   arriba que puedes trepar". Si un hijo no puede trepar por encima de
   ti, la arista o el nodo son críticos. Explicarlo así suma puntos.

⚠️ Cuidado con la raíz: la regla \verb|low[v] >= disc[u]| falla para la
   raíz del DFS. La raíz es articulación solo si tiene 2 o más hijos en
   el árbol DFS.

🔑 Aristas múltiples: si hay dos aristas paralelas entre u y v, ninguna
   es puente (hay ruta alterna). Trátalo con cuidado si el enunciado las
   permite.

\end{dsaentrevistas}

\begin{dsaresumen}

• Puente: arista que, al quitarla, desconecta el grafo.
• Punto de articulación: vértice que, al quitarlo, lo desconecta.
• Tarjan: un DFS con disc[] (descubrimiento) y low[] (mínimo alcanzable).
• Puente si low[v] > disc[u]; articulación si low[v] >= disc[u].
• Raíz es articulación solo con 2+ hijos en el DFS.
• Todo en O(V + E).
════════════════════════════════════════════════════════════════

\end{dsaresumen}
